\section{Related Work}
\label{sec:related}

Our work sits at the intersection of active perception, deformable-object manipulation, online system identification, and cross-material simulation benchmarks. The architectural pieces of \PTA{}---a fixed open-loop probing segment, an MLP encoder of the probing trace, and a context-conditioned residual policy---are deliberately controlled and individually well precedented. We position \PTA{} not as a new inference architecture but as an instantiation of active in-episode material identification used to isolate when probing is beneficial or costly for deformable-material manipulation. Below we review each adjacent line of work, emphasizing what \PTA{} \emph{infers}, what it \emph{conditions on}, and what it is \emph{evaluated against}, rather than claiming algorithmic novelty over each.

\paragraph{Active perception for manipulation.}
Robots have long been studied as embodied inference agents that gather information by acting on the world. Early work on active vision~\cite{bajcsy1988active} and exploratory procedures for haptic object recognition~\cite{lederman1987hand} established the principle that specific motions yield specific physical information. In modern learning-based manipulation, tactile exploration has been used to estimate local object geometry and contact properties~\cite{calandra2018more,xu2013tactile}, to inform regrasping policies~\cite{bauza2020tactile}, and to estimate object pose through pushing~\cite{yu2016more,bauza2018data}. Recent interactive perception systems actively choose visuo-tactile pushing or pulling actions to infer physical object properties~\cite{dutta2025predictive}, and tactile-informed dynamics models use contact observations for dense non-prehensile manipulation~\cite{ai2024robopack}. These works motivate probing as a physical measurement operation, but many focus on rigid objects, local contact properties, or closed-loop controllers where the measurement action's downstream damage is not isolated. \PTA{} extends this line by targeting continuous, bulk physical properties of \emph{deformable} media (stiffness, cohesion, plasticity) while treating the probe itself as an intervention whose disturbance may or may not relax before the downstream push.

\paragraph{Deformable and granular manipulation.}
Manipulating soft and flowable materials has become a focus area as differentiable and particle-based physics engines have matured. Environments such as PlasticineLab~\cite{huang2022plasticinelab} and DiffSkill~\cite{lin2022diffskill} study skill learning on plasticine-like media via differentiable MPM simulation, and RoboCook-style systems~\cite{shi2023robocook,shi2022robocraft} study tool use on dough-like materials. Recent work also adapts graph-based dynamics to material properties~\cite{zhang2024adaptigraph}, identifies elastoplastic parameters from robot interaction~\cite{yang2025differentiable}, and scales granular robotic-task simulation~\cite{millard2023granulargym}. A common feature of these works is that training and evaluation use the same material family, learn or calibrate an explicit dynamics model, or treat material variation via blanket domain randomization. Our benchmark instead organizes evaluation around a \emph{cross-material} split structure---training on the hardest granular material and evaluating zero-shot on cohesive and elastoplastic media---so that the role and cost of material-specific probing can be isolated.

\paragraph{Online system identification and adaptation.}
Adapting policies to unseen dynamics has been approached through meta-learning of dynamics and policies~\cite{finn2017model,nagabandi2019learning}, context-conditioned policies that infer a latent task embedding from recent history~\cite{rakelly2019efficient,zintgraf2021varibad}, and rapid online model adaptation~\cite{nagabandi2018deep,clavera2019modelbased}. Robotic sim-to-real pipelines have combined domain randomization with online adaptation of a low-dimensional parameter vector~\cite{yu2017preparing,peng2018sim2real}. Active system-identification methods go further by optimizing exploratory behavior to collect data for simulator or parameter refinement~\cite{memmel2024asid}. These approaches typically assume that the information needed to adapt is present in the course of normal task execution, or that the identification signal's physical cost is secondary to the parameter estimate. \PTA{} instead \emph{budgets} a dedicated, short probing phase at the start of each episode whose only purpose is to elicit material response. The conditioning latent and the residual policy are otherwise standard. The contribution we attach to this design choice is empirical, not architectural: we use it as a controlled lens to ask when separating ``what material am I holding?'' from ``how do I act?'' helps versus hurts, and our results show that the answer depends on whether probe-induced perturbations relax on the task timescale.

\paragraph{Multi-physics simulation and cross-material benchmarks.}
Recent simulators have made multi-material robotics research practical. Genesis~\cite{genesis2024} provides a unified framework covering rigid body, MPM, SPH, and FEM physics. Isaac Gym~\cite{makoviychuk2021isaac} and Brax~\cite{freeman2021brax} are widely used for GPU-accelerated rigid-body RL but do not natively support deformable media. Benchmarks such as SoftGym~\cite{lin2021softgym}, PlasticineLab~\cite{huang2022plasticinelab}, and GranularGym~\cite{millard2023granulargym} evaluate policies on deformable or granular tasks but typically fix the material family per task. Our benchmark complements these by pairing a single task definition with a deliberately discriminative material sweep (55pp of baseline variation across families), along with parameter-shifted in-family OOD splits and a standardized evaluator. The goal is a narrow, reproducible testbed in which the specific contribution of material-adaptive components---as opposed to task-specific engineering---can be read off the numbers.
